\documentclass[letterpaper]{article}
\usepackage[submission]{aaai2027}
\usepackage[hyphens]{url}
\usepackage{graphicx}
\urlstyle{rm}
\def\UrlFont{\rm}
\usepackage{natbib}
\usepackage{caption}
\usepackage{subcaption}
\frenchspacing
\usepackage{booktabs}
\usepackage{multirow}

% Additional packages (double-check whether they are allowed in AAAI 2027)
\usepackage{amsmath}
\usepackage{amssymb}
\usepackage{enumitem}
\setlist[itemize]{itemsep=0pt,parsep=0pt,topsep=2pt,leftmargin=*}
\setlist[enumerate]{itemsep=0pt,parsep=0pt,topsep=2pt,leftmargin=*}
\usepackage{multicol}
\usepackage{listings}
\usepackage[most]{tcolorbox}
\tcbset{
  colback=gray!5!white,
  colframe=gray!75!black,
  fonttitle=\bfseries,
  arc=3mm,
  boxrule=0.5mm,
  left=2mm,right=2mm,top=1mm,bottom=1mm,
  breakable
}
\lstset{
  basicstyle=\small\ttfamily,
  breaklines=true,
  breakatwhitespace=true,
  columns=flexible,
  keepspaces=true,
  showstringspaces=false
}

% Tighten float and caption spacing to reduce page count
\setlength{\textfloatsep}{6pt plus 1pt minus 1pt}
\setlength{\intextsep}{6pt plus 1pt minus 1pt}
\setlength{\floatsep}{6pt plus 1pt minus 1pt}
\setlength{\dbltextfloatsep}{6pt plus 1pt minus 1pt}
\setlength{\dblfloatsep}{6pt plus 1pt minus 1pt}
\setlength{\abovecaptionskip}{4pt}
\setlength{\belowcaptionskip}{2pt}
\setlength{\abovedisplayskip}{4pt plus 1pt minus 1pt}
\setlength{\belowdisplayskip}{4pt plus 1pt minus 1pt}

\pdfinfo{
/TemplateVersion (2027.1)
}

\setcounter{secnumdepth}{0}

% \title{Beyond FID: An Agentic Evaluation System for Class-to-Image Generative Models via Multi-Expert Collaborative Assessment}
% \title{Beyond FID: An Agentic Multi-Expert System for Interpretable Per-Image Evaluation of c2i Generative Models}
\title{What Sub-One FID Misses?\\
A Per-Image Evaluation of Alignment and Artifact for Image Generation}

\author{Anonymous Author}
\affiliations{
    Anonymous Institution\\
    \texttt{anonymous@example.com}
}

\newcommand{\llmname}{Qwen3.6-Plus}

\begin{document}

\maketitle

\begin{abstract}
The Fr\'echet Inception Distance (FID) has long been the de facto metric for class-to-image (c2i) evaluation, yet it reduces an entire dataset to a single scalar with no per-image interpretability and relies on a Gaussian assumption empirically find to be violated. We propose THEMIS, an agentic evaluation framework that scores every generated image on two disentangled dimensions---Alignment (class conformance) and Authenticity (structural integrity)---and reports set-level Diversity via Vendi Ratio. THEMIS employs a four-stage pipeline: a Router plans the evaluation, a Judge reviews the plan through multi-round feedback, seven specialized vision experts execute the plan, and a Reflector synthesizes all evidence provided by experts into final scores with self-reflection and few-shot human calibration. This agentic design yields per-image diagnostic reports that pinpoint specific defects, enabling targeted model improvement. Across five c2i models, THEMIS achieves perfect ranking agreement with human judgment ($\rho = 1.0$) and system stability matching human inter-annotator reproducibility (ICC $\geq 0.70$), outperforming FID in both rank correlation and score-level fidelity.
\end{abstract}

\section{Introduction}

Class-to-image (c2i) generative models have advanced rapidly, producing images that are increasingly difficult to distinguish from real photographs. As models mature, the ability to precisely evaluate their quality---not just to rank them, but to diagnose where and why they fail---becomes essential for guiding further improvement. Yet the field still relies on a single metric that was designed a decade ago for a fundamentally different purpose.

The Fr\'echet Inception Distance (FID)~\citep{heusel2017gans} measures the distance between Gaussian-modeled feature distributions of real and generated sets. While effective as a coarse ranking signal in earlier eras, FID has two structural limitations that are becoming critical. First, it assumes that Inception features follow a unimodal Gaussian---an assumption we empirically demonstrate to be violated for fine-grained ImageNet categories, where varied poses, viewpoints, and backgrounds create multi-cluster feature spaces. Second, and more fundamentally, FID collapses an entire dataset into a single scalar, conflating class conformance with visual quality and offering no per-image diagnosis. A model that generates a flawless dog with a melted limb receives the same scalar signal as one that produces a structurally intact but semantically wrong image. As modern c2i models cluster around similar FID scores (e.g., several current models report FID $\leq 1.0$), this opacity makes it impossible to identify failure modes or guide targeted improvements.

Existing alternatives each address only part of the problem. Distribution-level metrics such as Clean FID~\citep{parmar2021cleanfid} and CMMD~\citep{jayasumana2024rethinking} correct statistical biases but remain set-level scalars with no per-image granularity. LLM-as-a-judge approaches~\citep{zheng2023judging} can produce per-image assessments but lack the specialized vision capabilities needed to detect structural defects such as anatomical anomalies or subtle artifacts. What is needed is an evaluation system that combines the reliability of statistical metrics with the interpretability of per-image diagnosis, grounded in specialized visual evidence.

We propose \textbf{THEMIS}, an agentic evaluation framework that addresses these limitations through three design principles: \textbf{per-image interpretability}, \textbf{agentic reliability}, and \textbf{dimensional disentanglement}. Rather than reducing evaluation to a distribution-level statistic, THEMIS scores every generated image individually on two disentangled dimensions---Alignment and Authenticity---and supplements these with set-level Diversity via Vendi Ratio. A four-stage pipeline realizes these principles: (1)~a \emph{Router} analyzes each image with taxonomy priors to formulate an evaluation plan; (2)~a \emph{Judge} reviews and iteratively refines the plan through multi-round feedback; (3)~seven specialized \emph{expert models} (classification, detection, pose estimation, etc.) execute the plan to provide structural evidence that the LLM alone cannot assess; and (4)~a \emph{Reflector} synthesizes all evidence into final scores via self-reflection, calibrated against few-shot human-annotated reference images. This agentic design---combining multi-agent collaboration, expert grounding, iterative feedback, and self-reflection---ensures that scoring is reliable and interpretable, producing per-image diagnostic reports that pinpoint specific defects.

Our contributions are as follows:
\begin{enumerate}
\item We propose THEMIS, an agentic framework for c2i evaluation that provides per-image interpretable scores on two disentangled dimensions (Alignment and Authenticity) plus set-level Diversity, comparing FID's single scalar.
\item We demonstrate that THEMIS achieves perfect ranking agreement with human judgment ($\rho = 1.0$) across five c2i models, substantially outperforming FID ($\rho = 0.667$), and that system stability matches human inter-annotator reproducibility (ICC $\geq 0.70$).
\item We establish a c2i evaluation benchmark comprising five state-of-the-art models evaluated on 500 images per model with paired human annotations, and empirically validate that 5 images per class suffice for stable and reproducible evaluation, providing a standardized protocol for future c2i assessment.
\end{enumerate}

Figure~\ref{fig:crossmodel_scatter} visualizes system-human agreement across five c2i models (min-max normalized to [0,1]). THEMIS (a) achieves NRMSE$=$0.046, closely tracking human judgments; FID (b) yields NRMSE$=$0.301, with iMF-XL-FD and JiT-H-FD collapsing to the same point despite different human scores.

\begin{figure}[t]
\centering
\includegraphics[width=\linewidth]{Figures/crossmodel_scatter.png}
\caption{System-Human agreement across five c2i models. Scores are min-max normalized to [0,1].}
\label{fig:crossmodel_scatter}
\end{figure}

\input{section/RelatedWork}

\input{section/Limitations_of_FID}

\input{section/Methodology}

\input{section/Experiment}

\section{Conclusion}

We presented THEMIS, an agentic evaluation framework for class-to-image generative models focusing not only on taxonomy class alignment but also on complete image authenticity. Its four-stage pipeline---combining multi-agent collaboration, specialized vision experts, multi-round feedback, and self-reflection---produces disentangled Alignment and Authenticity scores that align closely with human judgment while providing interpretable, per-image evidence beyond FID's single scalar. Though the agentic pipeline incurs higher cost than distribution-level metrics, our analysis shows that it is worthy for interpretability and system-human consistency. Our experiments also prove that 500 images are sufficient for stable and reproducible evaluation.

\bibliography{themis}

\input{section/Appendix}

\end{document}
